\subsection{Problem S2-VD: explicit radial domains}

For \(0<L<U<1\) and \(0\le x\le1\), define the initial-component endpoint
\[
 \mathsf E(L,U;x)=
 \begin{cases}
 x,&x<L,\\
 U,&L\le x\le U,\\
 x,&U<x,
 \end{cases}
\]
and the complementary residual
\[
 \mathcal R_{[L,U]}(x)=1-\mathsf E(L,U;x).
\]
This is an explicit three-cell function; no supremum or placement operation
occurs in the definition.

For \(t>0\), define
\[
 \Delta_t=\sqrt{t^2+t+1}.
\]

\subsubsection{Adjacent domain}

Use independent variables
\[
 (r,a,d,p_0,q_0,t,p_1,q_1)\in\mathbb R^8
\]
and the signed quantities
\[
 w=1-r,\quad e=\sqrt{1-rw},\quad
 \eta=1-e,\quad k=\eta+a+d.
\]
Set
\[
 A_R=\mathcal R_{[k/r,w+d]}(q_0),
 \qquad
 A_L=\mathcal R_{[k/w,r+a]}(p_0),
\]
\[
 \lambda_+=\frac{t(1-q_1)}{t+1},
 \qquad
 u_+=\frac{\Delta_t-p_1-tq_1-1}{t}.
\]

\begin{definition}[Adjacent Vd cells]
\label{def:s2-vd-adjacent-domain}
Let \(\Omega_{\rm VD}^{\rm adj,+}\) be the set of the displayed tuples
satisfying
\[
 (r,a,d)\in\mathcal D_2,
\]
\[
 0<p_0,q_0<1,\qquad
 p_0+q_0>1,\qquad
 p_0^2+p_0q_0+q_0^2\le1,
\]
\[
 A_L>0,\qquad A_R\ge0,\qquad A_R+A_L<\frac12,
\]
\[
 a+d<\frac1{24},\qquad
 a+d<\frac{\min\{r,w\}}6,
\]
\[
 t>0,\qquad
 0<p_1,q_1<1,\qquad
 p_1\ge A_R,\qquad q_1\ge A_L,\qquad
 p_1+q_1<\frac12,
\]
and
\[
 0\le\lambda_+<u_+\le1.
\]
Let \(\Omega_{\rm VD}^{\rm adj,0}\) satisfy all the conditions above through
\(p_1+q_1<1/2\), and replace the final positive-interval condition by
\[
 u_+\le\lambda_+.
\]
The two sets respectively encode a positive first raw interval and no positive
first raw interval.
\end{definition}

Define
\[
 \Psi_{\rm VD}^{\rm adj,+}
 =\max\left\{d-\frac{A_L}{4},\ u_+-1+d\right\},
 \qquad
 \Psi_{\rm VD}^{\rm adj,0}
 =d-\frac{A_L}{4}.
\]

\begin{definition}[Problem S2-VD-adjacent]
\label{prob:s2-vd-adjacent}
The adjacent optimization statements are
\[
 \boxed{
 \Psi_{\rm VD}^{\rm adj,+}<0
 \quad\text{on }\Omega_{\rm VD}^{\rm adj,+},
 }
\]
and
\[
 \boxed{
 \Psi_{\rm VD}^{\rm adj,0}<0
 \quad\text{on }\Omega_{\rm VD}^{\rm adj,0}.
 }
\]
\end{definition}

The quantity \(u_+\) is now the explicit graph value
\((\Delta_t-p_1-tq_1-1)/t\); it replaces the former placement-defined far
endpoint.  The second term is the supported-interval target, and the first is
the quarter-margin target for the other possible bridge.

\subsubsection{Nonadjacent domain}

Use independent variables
\[
 (r,a,d,p_0,q_0,t,p_v,q_v)\in\mathbb R^8.
\]
Define
\[
 x=\frac{k}{w},\qquad y=\frac{k}{r},
\]
\[
 A_R=\mathcal R_{[y,w+d]}(q_0),
 \qquad
 A_L=\mathcal R_{[x,r+a]}(p_0),
\]
and
\[
 C_v=\frac{\Delta_t-p_v-tq_v}{t+1}.
\]
Put
\[
 \sigma=a+d.
\]

For \(j\in\{2,4\}\), define
\[
 d_2=d,\qquad d_4=a.
\]
For \(j=3\), use two order cells:
\[
 \mathcal O_{3,0}:\quad \frac ar\le\frac dw,\quad d_3=\frac ar,
\]
\[
 \mathcal O_{3,1}:\quad \frac dw<\frac ar,\quad d_3=\frac dw.
\]
Similarly split the minimum \(m=\min\{A_R,A_L\}\) into
\[
 \mathcal M_0:\quad A_R\le A_L,\quad m=A_R,
 \qquad
 \mathcal M_1:\quad A_L<A_R,\quad m=A_L.
\]

\begin{definition}[Nonadjacent Vd cells]
\label{def:s2-vd-nonadjacent-domain}
For
\[
 j\in\{2,4\},\qquad \mu\in\{0,1\},
\]
let \(\Omega_{\rm VD}^{j,\mu}\) be the set of tuples satisfying
\[
 (r,a,d)\in\mathcal D_2,
\]
\[
 0<p_0,q_0<1,\qquad
 p_0+q_0>1,\qquad
 p_0^2+p_0q_0+q_0^2\le1,
\]
\[
 A_R>0,\qquad A_L>0,\qquad A_R+A_L<\frac12,
\]
\[
 \sigma<\frac x2,\qquad \sigma<\frac y2,
\]
\[
 x^2+x(w+d)+(w+d)^2\le1,\qquad
 y^2+y(r+a)+(r+a)^2\le1,
\]
\[
 t>0,\qquad
 0<p_v,q_v<1,\qquad
 p_v\ge A_R,\qquad q_v\ge A_L,
\]
\[
 0\le C_v\le1,
\]
together with \(\mathcal M_\mu\).  For \(j=3\), additionally choose
\(\nu\in\{0,1\}\) and impose \(\mathcal O_{3,\nu}\); denote the resulting
cell by \(\Omega_{\rm VD}^{3,\mu,\nu}\).
\end{definition}

On every nonadjacent cell, define
\[
 \Psi_{\rm VD}^{\rm non}
 =\max\{\sigma-A_R,\ \sigma-A_L,\ d_j-m,\ C_v-1+m\}.
\]

\begin{definition}[Problem S2-VD-nonadjacent]
\label{prob:s2-vd-nonadjacent}
The nonadjacent optimization problem is
\[
 \boxed{
 \Psi_{\rm VD}^{\rm non}<0
 }
\]
on every cell
\[
 \Omega_{\rm VD}^{2,\mu},\quad
 \Omega_{\rm VD}^{4,\mu},\quad
 \Omega_{\rm VD}^{3,\mu,\nu}
 \qquad(\mu,\nu\in\{0,1\}).
\]
\end{definition}

\begin{definition}[Problem S2-VD]
\label{prob:s2-vd-terminal}
Problem S2-VD is the finite family consisting of
Problems~\ref{prob:s2-vd-adjacent} and
\ref{prob:s2-vd-nonadjacent}.
\end{definition}

All former interval-component suprema have been replaced by the explicit
three-cell function \(\mathsf E\), and all former radial placement suprema
have been replaced by the exact corner graph formulas for \(u_+\) and \(C_v\).
